% Prose for this counterexample.  Every region below is delimited by an
% environment that tools/build.py extracts; machine facts (ids, uid, dates,
% urls, enums) live in case.json.  See counterexamples/AGENTS.md.

\cxtitle{Variance-sensitive Matrix Spencer}

\begin{cxcredits}
\posedby{A.~S.~Bandeira \citeyearpar{Bandeira2016}}
\foundby{GPT-5.5 (Pro)}{TODO}
\formalizedby{E.~Akba\c{s} and S.~Sra \citeyearpar{AkbasSra2026}}
\auditedby{S.~Sra}
\contributedby{S.~Sra}
\end{cxcredits}

\begin{cxcontext}
For symmetric matrices $A_1,\ldots,A_n$ the \emph{discrepancy} is $\disc(A_1,\ldots,A_n)=\min_{x\in\{\pm1\}^n}\|\sum_{i=1}^nx_iA_i\|_{\rm op}$.  Spencer's theorem gives an $O(\sqrt n)$ bound in the scalar case, and the Matrix Spencer conjecture asks for the same $O(\sqrt n)$ bound for $n$ symmetric contractions of size $n\times n$.  The variance-sensitive strengthening below replaces the ambient scale $\sqrt n$ by the intrinsic variance parameter $\|\sum_{i=1}^nA_i^2\|_{\rm op}^{1/2}$, which never exceeds $\sqrt n$ and is often far smaller; it is the parameter that controls the norm of a random signing in matrix concentration, and Kyng, Luh and Song established the strengthened bound in the rank-one case.  The dimension regime is the whole point: the conjecture concerns $n$ matrices of size $n\times n$, so a family whose dimension grows faster than the number of summands says nothing about it.
\end{cxcontext}

% ---------------- result: variance-only-matrix-discrepancy ----------------

\begin{cxsource}{variance-only-matrix-discrepancy}
Remark 4.25 of A.~S.~Bandeira's problem collection \citeyearpar{Bandeira2016}, quoted here from its restatement as Conjecture 1.2 of E.~Akba\c{s} and S.~Sra \citeyearpar{AkbasSra2026}
\end{cxsource}

\begin{cxstatement}{variance-only-matrix-discrepancy}
There exists a universal constant $C$, such that for any choice of $n$ symmetric matrices $A_1,\ldots,A_n\in\R^{n\times n}$ satisfying $\|A_i\|_{\rm op}\le1$ there exists a coloring $x\in\{\pm1\}^n$, such that
\[\left\|\sum_{i=1}^nx_iA_i\right\|_{\rm op}\le C\left\|\sum_{i=1}^nA_i^2\right\|_{\rm op}^{1/2}.\]
\end{cxstatement}

\begin{cxsummary}{variance-only-matrix-discrepancy}
Variance-sensitive Matrix Spencer for $n$ contractions of size $n\times n$ \cite{Bandeira2016}
\end{cxsummary}

\begin{cxcertificate}{variance-only-matrix-discrepancy}
Exact $d=n$ diagonal family: $\disc=(m-1)/\sqrt m$ against variance $1+1/m$, a ratio of $\Theta(\sqrt{\log n})$.
\end{cxcertificate}

\begin{cxrefutation}
The disproof recorded here is not ours: the family below is Theorem A.1 of \citet{AkbasSra2026}, the first refutation of the conjecture, and this entry reproduces it.

Fix an integer $m\ge2$ and set $n=2^m$, so that the matrices below are $n\times n$ and there are $n$ of them: the dimension is not inflated relative to the number of summands.  Index coordinates by $s=(s_1,\ldots,s_m)\in S:=\{\pm1\}^m$, so $|S|=n$.  Write $p=(1,\ldots,1)$ and $q=(-1,\ldots,-1)$, choose any $U\subseteq S\setminus\{p,q\}$ with $|U|=m$, put $F=S\setminus U$, and fix a bijection $\pi:\{m+1,\ldots,n\}\to F$; note that $p,q\in F$.  Define diagonal matrices
\begin{equation}\label{eq:disc-family}
(A_i)_{s,s}=\frac{s_i}{\sqrt m}\quad(1\le i\le m),
\qquad
A_k=\frac1{\sqrt m}E_{\pi(k),\pi(k)}\quad(m<k\le n),
\end{equation}
where $E_{r,r}$ is the diagonal matrix with a single $1$ in coordinate $r$.  Each $\|A_i\|_{\rm op}=1/\sqrt m\le1$, so the family is admissible.

\begin{theorem}[\statusfalse: variance-sensitive Matrix Spencer]\label{thm:discrepancy}
For the family \eqref{eq:disc-family},
\[
\disc(A_1,\ldots,A_n)=\frac{m-1}{\sqrt m},
\qquad
\left\|\sum_{i=1}^nA_i^2\right\|_{\rm op}^{1/2}=\sqrt{1+\tfrac1m},
\]
so the ratio of the two sides equals $(m-1)/\sqrt{m+1}=\Theta(\sqrt{\log n})$.  No universal constant $C$ can therefore satisfy the variance-sensitive bound, already for $n$ matrices of size $n\times n$.
\end{theorem}
\begin{proof}
Since the matrices are diagonal, so is $\sum_iA_i^2$, and its $s$th entry is
\[
\sum_{i=1}^m\frac{s_i^2}m+\sum_{k=m+1}^n\frac1m\mathbf 1\{s=\pi(k)\}=1+\frac1m\mathbf 1\{s\in F\},
\]
whence $\|\sum_iA_i^2\|_{\rm op}=1+\frac1m$.

For the discrepancy, let $x\in\{\pm1\}^n$ be an arbitrary signing and $X=\sum_{i=1}^nx_iA_i$, so that
\[
X_{s,s}=\frac1{\sqrt m}\sum_{i=1}^mx_is_i+\frac1{\sqrt m}\mathbf 1\{s\in F\}x_{\pi^{-1}(s)}.
\]
Evaluating at the coordinate $s=(x_1,\ldots,x_m)\in S$ makes the first sum equal to $m$, so
\[
\|X\|_{\rm op}\ge|X_{s,s}|\ge\frac{m}{\sqrt m}-\frac1{\sqrt m}=\frac{m-1}{\sqrt m}.
\]
This bound is attained.  Take $x_1=\cdots=x_m=1$ and $x_{k_p}=-1$, $x_{k_q}=1$, where $k_p=\pi^{-1}(p)$ and $k_q=\pi^{-1}(q)$, with the remaining signs arbitrary.  Then $\sqrt mX_{s,s}=\sum_{j=1}^ms_j+\mathbf 1\{s\in F\}x_{\pi^{-1}(s)}$, which equals $m-1$ at $s=p$ and $-(m-1)$ at $s=q$.  For every other $s$ the vector is neither all-plus nor all-minus, so $|\sum_js_j|\le m-2$ and hence $|\sqrt mX_{s,s}|\le m-1$.  Thus $\|X\|_{\rm op}=(m-1)/\sqrt m$, and the infimum over signings is exactly that.

Dividing the two displayed quantities gives $(m-1)/\sqrt{m+1}$, which grows without bound as $m=\log_2n\to\infty$.
\end{proof}

\begin{remark}[Why the dimension regime matters]
An earlier version of this entry used $n$ matrices of size $2^n\times2^n$.  That is the trivial regime: once the dimension is exponential in the number of summands the failure says nothing about Matrix Spencer, whose whole content lies in $d$ comparable to $n$.  The family above is the informative one, with $d=n$.

The family is commutative, and the growth rate $\sqrt{\log n}$ is exactly what the sharp noncommutative Khintchine inequalities predict in the commutative case, so the logarithmic factor is genuine rather than an artifact of the construction.  This leaves open the possibility raised in \citet{AkbasSra2026} that the variance-sensitive bound survives for sufficiently noncommutative families.  The original Matrix Spencer conjecture is untouched: bounding $\|\sum_iA_i^2\|_{\rm op}^{1/2}\le\sqrt n$ recovers only the weaker $O(\sqrt n)$ target, which this family satisfies comfortably.
\end{remark}
\end{cxrefutation}
